% Part: first-order-logic
% Chapter: tableaux
% Section: soundness-identity

\documentclass[../../../include/open-logic-section]{subfiles}

\begin{document}

\olfileid{fol}{tab}{sid}

\olsection{Kasahihan mawa \usetoken{S}{identity}}

\begin{prop}
Yèn $\Gamma$ nduwèni !!{tableau} katutup nganggo aturan identitas,
ora ana struktur kang nyukupi kabèh formula mawa tandha ing $\Gamma$.
\end{prop}

\begin{proof}
Kita nganggo induksi ing cacah aplikasi aturan, kaya ing bukti
\olref[fol][tab][sou]{thm:tableau-soundness}. Kanggo aturan dhasar, anggo kasus-kasus ing
teorema iku. Kanggo saben aturan identitas, kita nuduhaké yèn aturan
iku njaga kahanan bisa dicukupi ing cabang kang ditambahi.
Tetepna $\Gamma$ minangka himpunan !!{formula}s mawa tandha ing
cabang lan $\Struct{M}$ minangka !!a{structure} kang nyukupi~$\Gamma$.

Upamana cabang diambakake nganggo~$\eq$, yaiku kanthi nambahake
!!{formula} mawa tandha~$\sFmla{\True}{\eq[t][t]}$. Mesthi wae,
$\Sat{M}{\eq[t][t]}$, mula $\Struct{M}$ uga nyukupi $\Gamma \cup
\{\sFmla{\True}{\eq[t][t]}\}$.

Yèn cabang diambakaké nganggo $\TRule{\True}{\eq}$, aturan nambahaké
!!{formula} mawa tandha~$\sFmla{\True}{!A(t_2)}$. Kanggo
$\TRule{\False}{\eq}$, tandha premis formula lan tandha kesimpulan
padha-padha $\False$. Ing loro kasus, $\Gamma$ ngemot
$\sFmla{\True}{\eq[t_1][t_2]}$ lan $\sFmla{S}{!A(t_1)}$.
Term $t_1$ lan $t_2$ katutup, lan amarga padha nilainé ing
$\Struct{M}$, tetepa $d=\Value{t_1}{M}=\Value{t_2}{M}$.
Tetepna pangaji $s$ kanthi $s(x)=d$. Miturut
\olref[syn][ext]{prop:ext-formulas}, amarga ngganti pangaji $x$ nganggo
$d$ tetep mènèhi $s$, diduwèni
\[
  \Sat{M}{!A(t_1)}[s] \quad\text{yèn lan mung yèn}\quad
  \Sat{M}{!A(x)}[s] \quad\text{yèn lan mung yèn}\quad \Sat{M}{!A(t_2)}[s].
\]
Miturut \olref[syn][ass]{prop:sentence-sat-true}, loro instansi iku
ukara lan ora gumantung marang pangaji, mula
$\Sat{M}{!A(t_1)}$ yen lan mung yen $\Sat{M}{!A(t_2)}$.
Mula $\Struct{M}$ nyukupi $\sFmla{S}{!A(t_2)}$ kanggo tandha $S$
kang padha karo tandha premis formula.
\end{proof}

\end{document}
